\section{Mô tả chức năng và dữ liệu}\label{sec:3}

Mỗi descriptor tính $C = A \times B$ trên INT8, tích lũy INT32. Khi \sig{POST_EN} = 0, C (INT32) được
ghi ra DRAM. Khi \sig{POST_EN} = 1, VPU biến C thành
$Y = \text{kẹp}(\text{ReLU}(\text{requantize}(C + bias)))$ ở INT8 và Y được ghi ra DRAM. Mọi phép
tính là số nguyên và kết quả phải khớp từng bit với mô hình tham chiếu ở mục~\ref{sec:3.5}.

\subsection{Phép toán}\label{sec:3.1}

\paragraph{Phép toán}\label{par:3-1}
\begin{itemize}
\item $C[i][j] = \sum_{k=0}^{K-1} A[i][k]\cdot B[k][j]$, với 0 ≤ i < M, 0 ≤ j < N
\item A, B là INT8 bù 2. C và bias là INT32 bù 2. Y là INT8 bù 2. Tích INT8 × INT8 tính đủ 16 bit có dấu rồi mở rộng dấu lên 32 bit; cộng dồn trên INT32, wrap mod $2^{32}$ nếu tràn
\item Khi \sig{POST_EN} = 1, với mỗi phần tử: \newline
\hspace*{1em}$acc = C[i][j] + bias[j]$ \hfill (INT32, wrap mod $2^{32}$) \newline
\hspace*{1em}$prod = acc \cdot M_0$ \hfill (tích có dấu 64 bit) \newline
\hspace*{1em}$y = (prod + 2^{30+sh}) \gg (31+sh)$ \hfill (dịch số học, 64 bit) \newline
\hspace*{1em}nếu \sig{RELU_EN} = 1: $y = \max(y, 0)$ \newline
\hspace*{1em}$Y[i][j] = \min(\max(y, -128), 127)$
\item $M_0$ (\sig{QM0} của descriptor) là INT32 có dấu, mọi giá trị đều hợp lệ. $sh$ (\sig{QSHIFT}) là số không dấu 0 … 31. Một cặp ($M_0$, $sh$) dùng cho một descriptor. bias là vector N phần tử INT32, cộng theo cột
\item Kết quả ghi đè vùng đích trong DRAM; accelerator không đọc giá trị cũ của vùng đích
\item Chỉ hỗ trợ nhân ma trận và hậu xử lý đã mô tả; không có $C \mathrel{+}= A \times B$, softmax hoặc pooling
\end{itemize}

\begin{notebox}[Hệ quả]
$|a\cdot b| ≤ 2^{14}$ và K ≤ 32 nên $|C| ≤ 2^{19}$: C không bao giờ tràn. $|prod| ≤ 2^{62}$ và
$2^{30+sh} ≤ 2^{61}$ nên phép tính hậu xử lý 64 bit không bao giờ tràn. Làm tròn là nửa lên
(về phía $+\infty$). Phần mềm nên chọn $M_0$ trong $[2^{30}, 2^{31})$ để giữ độ chính xác.
\end{notebox}

\subsection{Kiểu dữ liệu, dấu và tràn}\label{sec:3.2}

\begin{spectable}{|L{0.22}|L{0.16}|L{0.62}|}{Kiểu dữ liệu và quy tắc tràn}{tab:3-2}{\textbf{Đại lượng} & \textbf{Kiểu} & \textbf{Quy tắc}}
A, B & INT8 có dấu & Bù 2. Không có dạng không dấu \\\hline

$a \cdot b$ & 16 bit có dấu & Tính đủ 16 bit rồi \textbf{mở rộng dấu} lên 32 bit trước khi cộng \\\hline

C (tích lũy) & INT32 có dấu & \textbf{Wrap} mod $2^{32}$, không saturate. Trên miền hợp lệ không bao giờ tràn ($|C| ≤ 2^{19}$) \\\hline

bias & INT32 có dấu & $acc = C + bias[j]$ trên INT32, \textbf{wrap} mod $2^{32}$ \\\hline

$M_0$ (\sig{QM0}) & INT32 có dấu & Mọi giá trị hợp lệ, kể cả $-2^{31}$ \\\hline

$sh$ (\sig{QSHIFT}) & 5 bit không dấu & Miền 0 … 31 \\\hline

$prod$ & 64 bit có dấu & $acc \cdot M_0$; không tràn vì $|prod| ≤ 2^{62}$ \\\hline

$y$ trước kẹp & 64 bit có dấu & Dịch \textbf{số học}; làm tròn nửa lên (về $+\infty$) \\\hline

Y & INT8 có dấu & \textbf{Saturate} về $[-128, 127]$. Đây là chỗ \textbf{duy nhất} trong thiết kế có saturation \\\hline

\end{spectable}

\begin{notebox}[Quyết định baseline D-27]
Hai nguồn nêu rải rác và không hoàn toàn trùng cách diễn đạt về dấu và tràn; baseline gom thành một
quy tắc duy nhất trong bảng~\ref{tab:3-2}: \textbf{mọi phép cộng số nguyên wrap, chỉ bước tạo Y
saturate}. Hiện thực requantize trong VPU dùng hai bước tương đương biểu thức ở
mục~\ref{sec:3.1}: $t = prod \ggg (30+sh)$ giữ đủ \textbf{34 bit}, rồi $y = (t+1) \ggg 1$ giữ đủ
\textbf{33 bit}, rồi kẹp với $lo = \text{\sig{RELU_EN}} \mathbin{?} 0 : -128$.
Bắt buộc giữ đủ 34 và 33 bit: cắt $t$ về 32 bit cho kết quả \textbf{sai} tại
$acc = M_0 = -2^{31}$. Đây là một lỗi hiện thực đã được nhận diện trong nguồn và baseline nâng lên
thành yêu cầu bắt buộc, kèm điểm kiểm chứng tương ứng trong spec của VPU.
\end{notebox}

\subsection{Kích thước ma trận và căn chỉnh}\label{sec:3.3}

\label{tab:3-3}
Miền kích thước của một descriptor: $1 ≤ M ≤ 32$; $N, K \in \{4, 8, \dots, 32\}$.

\begin{spectable}{|L{0.2}|L{0.8}|}{Quy tắc kích thước và căn chỉnh}{tab:3-4}{\textbf{Hạng mục} & \textbf{Quy tắc}}
M & Bất kỳ giá trị 1 … 32. Không cần là bội của 4, không cần là bội của P \\\hline

N, K & Bội của 4 trong miền 4 … 32. Ma trận có N hoặc K không chia hết cho 4 phải được \textbf{phần mềm} đệm 0 lên bội của 4; cột thừa của kết quả do phần mềm bỏ qua \\\hline

Không chia hết cho P = 8 & Xử lý \textbf{trong phần cứng}, không cần đệm: $TM = \lceil M/8 \rceil$, $TN = \lceil N/8 \rceil$; tile biên dùng \sig{row_vld[r]} = (i0 + r < M) chặn hàng thừa và \sig{o_row_mask} (8 bit) chặn cột thừa \\\hline

K không chia hết cho 8 & Không cần xử lý gì: K là chiều tích lũy, bộ tuần tự tile phát đúng K bước rồi chèn bubble tới $I = \max(K, 8)$ \\\hline

Địa chỉ cơ sở & Mọi địa chỉ cơ sở được dùng phải căn 4 byte. \sig{DESC_PTR} và mọi \sig{NEXT_PTR} được dùng phải căn 64 byte \\\hline

Biên word & Vì N và K là bội của 4, mọi hàng bắt đầu và kết thúc đúng biên word, nên \sig{m_axi_wstrb} luôn bằng 4'b1111 \\\hline

Dung lượng bank & A, B, bias và C của một descriptor nằm trọn trong một bank của buffer tương ứng: bank A và B sâu 256 word, C 1024 word, BIAS 32 word, đủ cho M, N, K ≤ 32 \\\hline

\end{spectable}

\begin{notebox}[Quyết định baseline D-26]
Baseline chốt ranh giới trách nhiệm: \textbf{phần mềm} đệm tới bội của 4, \textbf{phần cứng} xử lý
phần dư so với P = 8. Lý do: đệm tới bội 4 là điều kiện đủ để mọi hàng DRAM trùng biên word 32 bit,
nhờ đó DMA không bao giờ cần \sig{wstrb} từng phần; còn buộc phần mềm đệm tới bội 8 sẽ đòi thêm bộ
nhớ và thêm bước xử lý mà phần cứng vốn đã có đủ cơ chế (\sig{row_vld} và \sig{o_row_mask}) để
chặn. Hai cơ chế chặn khác nhau: \textbf{hàng} thừa có \sig{row_vld[r]} = 0 nên PE của hàng đó không
cập nhật \sig{acc} và không phát \sig{trig}, tức hàng đó không bao giờ được xả; \textbf{cột} thừa
vẫn được PE tính bình thường nhưng kết quả không được ghi, vì \sig{o_row_mask} đi tiếp thành
\sig{i_c_wr_mask} rồi thành byte/word enable của bank C.
\end{notebox}

\subsection{Bố trí dữ liệu trong DRAM}\label{sec:3.4}

\begin{spectable}{|L{0.14}|L{0.46}|L{0.4}|}{Địa chỉ byte của phần tử}{tab:3-5}{\textbf{Ma trận} & \textbf{Địa chỉ khi stride = 0} & \textbf{Kích thước vùng}}
A & \sig{A_ADDR} + i·K + k & M·K byte \\\hline

B & \sig{B_ADDR} + k·N + j & K·N byte \\\hline

bias & \sig{BIAS_ADDR} + 4j & 4·N byte \\\hline

C (INT32) & \sig{C_ADDR} + 4·(i·N + j) & 4·M·N byte \\\hline

Y (INT8) & \sig{C_ADDR} + i·N + j & M·N byte \\\hline

\end{spectable}

\paragraph{Bố trí dữ liệu}\label{par:3-6}
\begin{itemize}
\item Địa chỉ trong bảng là lúc stride bằng 0 (các hàng liền nhau). Khi \sig{LDA}, \sig{LDB}, \sig{LDC} khác 0 thì dùng công thức ở mục~\ref{sec:5.2}
\item Little-endian: byte tại địa chỉ a nằm ở bit [8·(a mod 4)+7 : 8·(a mod 4)] của word 32 bit; INT32 là little-endian
\item Mọi địa chỉ cơ sở căn 4 byte. Vì N, K là bội của 4, mọi hàng bắt đầu và kết thúc đúng biên word, nên \sig{m_axi_wstrb} luôn bằng 4'b1111
\item Vùng đích của một descriptor không được chồng lên A, B hoặc bias của descriptor đó; A, B, bias được phép chồng nhau. Vùng đích của mọi descriptor không được chồng lên vùng descriptor. Vi phạm → kết quả không xác định. Phần cứng không kiểm tra
\item Phần cứng không kiểm tra vùng nhớ có nằm trong DRAM không; bus trả lỗi thì tác vụ kết thúc với ERROR (mục~\ref{sec:4.4})
\end{itemize}

\begin{notebox}[Quyết định baseline D-24]
Phần cứng \textbf{không} kiểm tra giới hạn địa chỉ DRAM và \textbf{không} kiểm tra chồng lấn vùng
nhớ. Lý do: kiểm chồng lấn cần biết toàn bộ bản đồ bộ nhớ của SoC và cần so sánh mọi cặp vùng của
cả chuỗi, không khả thi trong accelerator. Phần cứng chỉ kiểm những điều kiện cục bộ của descriptor
ở mục~\ref{sec:5.2} (căn chỉnh, miền M/N/K, stride, bit dự trữ). Tác động: vi phạm chồng lấn cho
kết quả không xác định mà không báo lỗi; truy cập ngoài vùng DRAM sẽ do bus phát hiện và trở thành
\sig{ERR_CODE} = 2, 3 hoặc 4.
\end{notebox}

\subsection{Mô hình tham chiếu}\label{sec:3.5}

\begin{Verbatim}[frame=leftline,framerule=2pt,rulecolor=\color{brand},xleftmargin=6pt,samepage=true]
import numpy as np

def gemm(A, B):                      # A: int8 MxK, B: int8 KxN
    return (A.astype(np.int64) @ B.astype(np.int64)).astype(np.int32)

def post(C, bias, M0, sh, relu):     # C: int32 MxN, bias: int32 N
    acc  = (C.astype(np.int64) + bias).astype(np.int32).astype(np.int64)
    prod = acc * np.int64(M0)
    y    = (prod + (np.int64(1) << (30 + sh))) >> (31 + sh)
    if relu:
        y = np.maximum(y, 0)
    return np.clip(y, -128, 127).astype(np.int8)
\end{Verbatim}

Nội dung DRAM sau một tác vụ phải khớp từng bit với mô hình trên, áp dụng lần lượt cho từng
descriptor theo đúng thứ tự chuỗi, bất kể các pha đã chạy chồng nhau đến mức nào.

\subsection{Ma trận lớn hơn 32}\label{sec:3.6}

Một descriptor giới hạn M, N, K ≤ 32. Lớp lớn hơn được phần mềm chia thành nhiều descriptor:

\begin{spectable}{|L{0.14}|L{0.86}|}{Chia lớp lớn hơn 32}{tab:3-7}{\textbf{Chiều} & \textbf{Cách chia}}
M & Chia thành các dải 32 hàng. Mỗi dải là một descriptor độc lập; A, C dùng địa chỉ cơ sở lệch theo hàng và stride \sig{LDA}, \sig{LDC} bằng độ dài hàng thật \\\hline

N & Chia thành các dải cột ≤ 32. Dải bắt đầu ở cột $j_0$ dùng đoạn bias tương ứng (\sig{BIAS_ADDR} lệch $4j_0$ byte, tức 32 word = 128 byte mỗi dải 32 cột), \sig{B_ADDR} lệch $j_0$ byte, \sig{C_ADDR} lệch $4j_0$ byte (C INT32) hoặc $j_0$ byte (Y INT8); \sig{LDB}, \sig{LDC} bằng độ dài hàng thật; \sig{REUSE_B} không dùng được giữa các dải cột khác nhau \\\hline

K & Chia thành các đoạn ≤ 32 với \sig{POST_EN} = 0, ghi ra các vùng C INT32 riêng; CPU cộng dồn các vùng đó rồi tự thực hiện bias, requantize, ReLU và kẹp. Không có chế độ $C \mathrel{+}= A \times B$ trong phần cứng \\\hline

\end{spectable}

Nội dung mục này đến từ \sig{spec_2} (REQ-AP-001); \sig{spec_1} không có phần tương ứng. Baseline
giữ lại vì nó cần thiết để phần mềm dùng được accelerator cho lớp mạng thật, và nó không mâu thuẫn
với bất kỳ quyết định nào ở trên.

\subsection{Thông lượng}\label{sec:3.7}

\begin{spectable}{|L{0.42}|L{0.26}|L{0.32}|}{Thông lượng baseline, M = N = K = 32, \sigc{POST_EN} = 0, stride = 0}{tab:3-8}{\textbf{Hạng mục} & \textbf{Khối lượng (word, beat hoặc chu kỳ)} & \textbf{Thời gian ở 100 MHz}}
FETCH & 16 word, 1 burst 16 beat & 0,16 µs \\\hline

LOAD\_A & 32 hàng × 8 word = 256 word & 2,56 µs \\\hline

LOAD\_B & 32 hàng × 8 word = 256 word & 2,56 µs \\\hline

STORE\_C (INT32) & 32 hàng × 32 word = 1024 word & 10,24 µs \\\hline

Tổng beat AXI & 1552 beat & 15,52 µs \\\hline

COMPUTE & 534 chu kỳ & 5,34 µs \\\hline

Đỉnh tính toán của array & 64 MAC/chu kỳ & 6,4 GMAC/s (12,8 GOP/s) \\\hline

\end{spectable}

\begin{notebox}[Quyết định baseline D-25]
Hai nguồn đưa các số liệu thành phần nhưng không nguồn nào tổng hợp thành bức tranh thông lượng;
baseline tính và chốt bảng~\ref{tab:3-8} từ các số liệu đó. Kết luận dùng cho thiết kế:
với M = N = K = 32 và \sig{POST_EN} = 0, tác vụ \textbf{bị giới hạn bởi băng thông AXI 32 bit}
chứ không bởi array (534 chu kỳ ≈ 5,34 µs): một descriptor chạy đơn lẻ cần 1552 beat ≈ 15,52 µs;
khi chạy chồng, kênh đọc (528 beat) và kênh ghi (1024 beat) hoạt động song song nên khâu chậm nhất
là STORE\_C (≈ 10,24 µs). Vì vậy scratchpad ping-pong che
được toàn bộ thời gian COMPUTE, và thời gian mỗi descriptor tiến tới $\max(\text{nạp}, \text{tính},
\text{ghi}) \approx$ thời gian STORE\_C. Hệ quả thiết kế: không cần làm array nhanh hơn; cải thiện
thông lượng phải đến từ độ rộng bus AXI hoặc từ \sig{POST_EN} = 1 (STORE\_C giảm bốn lần, còn
256 word). Với \sig{POST_EN} = 1, tổng beat là 16 + 32 + 256 + 256 + 256 = 816, tức 8,16 µs, và
COMPUTE (5,34 µs) bắt đầu trở thành đáng kể.
\end{notebox}
